
\chapter{Conclusiones}

El desarrollo del sistema de alimentación automatizado para gallinas permitió alcanzar los objetivos planteados al inicio del proyecto. Se diseñó un sistema que integra correctamente cada uno de los módulos propuestos. Este diseño utiliza materiales adecuados como acero galvanizado, que proporcionaron resistencia y durabilidad al sistema, pensando en su uso con animales y resistencia a la corrosión.

Los resultados obtenidos en las pruebas mostraron que el sistema distribuye 3 kg de alimento a lo largo de 1.5 m aproximandamente, gracias al uso de sensores y motores controlados por la tarjeta ESP32 en conjunto con los circuitos integrados antes mencionados. Asimismo, la incorporación de una pantalla Nextion permite una interacción sencilla y efectiva para el usuario, mejorando la usabilidad del sistema y reduciendo la cantidad de periféricos.

Tras la integración del sistema se detectan áreas de oportunidad en la selección de materiales y geometrías del sistema. Los elementos mecánicos meaquinados cuentan con tolerancias de hasta 1 mm provocando vibraciones. 

\section{Recomendaciones y trabajo futuro}

Se prevé en un futuro diseño cambios en la geometría de la tolva fija, debido a que presentó problemas en la fabricación; en conjunto un rediseño en la estructura principal, reduciendo a PTR de $1\frac{1}{2}$ pulgadas o de $1$ pulgada, también se contempla aluminio como material para la estructura. 

Para la programación del módulo de alimentador se propone cambiar la \textit{MCU} por una FPGA que nos puede brindar una reducción en el uso de circuitos periféricos y mejor control en el hardware y menor consumo de corriente. 

%En el caso del módulo de monitoreo se propone la elaboración de un sistema RFID capaz de funcionar a $134.2 \text{ kHz}$ que pueda aprovechar correctamente las etiquetas especializadas.

El sistema contempla dos puntos importantes de un problema general, por lo que se propone la integración de sistemas capaces de mejorar la producción de huevo, tales como el control del ambiente; humedad y temperatura; o la implementación de sistemas capaces de monitorizar la salud de los animales; así mismo, se plantea la expansión para poblaciones más grandes de gallinas.



%%%%%%fin del archivo
\endinput 